\begin{textsample}{NEG Dim 1 – persona\_grok – Score -3.00 – t518\_grok.txt}  \label{ex:f1_neg_021}
Indonesia is once again shrouded in a smoky haze from raging fires, leading to severe health problems for millions and releasing massive amounts of CO2 into the atmosphere.
These fires are closely tied to \textbf{palm} oil production and rampant deforestation.
Greenpeace campaigner Annisa Rahmawati has called out major companies including Unilever, Mondelez, Nestle, and P\&G for their connections to nearly 10,000 fire hotspots in 2019 and over 200,000 hectares of burned land between 2015 and 2018, even though these companies have made commitments to no-deforestation policies.
The situation highlights failures in implementation, such as with Wilmar, which has not fully enforced its No Deforestation, No Peat, No Exploitation ( NDPE ) policy.
For instance, Wilmar initially suspended purchases from the problematic \textbf{supplier} GAMA/KPN but later resumed them.
Due to a lack of agreement on transparency, Greenpeace has withdrawn from collaborative \textbf{monitoring} efforts.
We are urging these companies to strictly enforce responsible sourcing practices or to abandon high-risk \textbf{commodities} altogether.

% matched lemmas: commodity, monitoring, palm, supplier
\end{textsample}
